% =============================================================================
% Pretrain-loss-weighting choice for neural-network density functionals.
% Cross-step comparison of unweighted vs. Becke-quadrature-integrated objectives.
%
% IMPORTANT BUILD ORDER:
%   1. Compile supplement.tex first (generates supplement.aux for cross-refs).
%   2. Then compile this main.tex.
%
% Local build:
%   cd reports_local/latex
%   mkdir -p build_supp build_main
%   # Step 1 — supplement
%   pdflatex -output-directory=build_supp supplement.tex
%   (cd build_supp && BIBINPUTS=.. bibtex supplement)
%   pdflatex -output-directory=build_supp supplement.tex
%   pdflatex -output-directory=build_supp supplement.tex
%   # Make supplement.aux visible to main (xr-hyper reads it)
%   cp build_supp/supplement.aux .
%   # Step 2 — main
%   pdflatex -output-directory=build_main main.tex
%   (cd build_main && BIBINPUTS=.. bibtex main)
%   pdflatex -output-directory=build_main main.tex
%   pdflatex -output-directory=build_main main.tex
%
% Overleaf:
%   - Add both main.tex and supplement.tex to the project.
%   - Set supplement.tex as the main document, compile (generates aux files).
%   - Set main.tex as the main document; xr-hyper resolves S: cross-refs
%     from supplement.aux automatically.
%
% Style: generic article class for portability; switch the documentclass to
% \documentclass[twocolumn,aps,prl]{revtex4-2} for PRL submission, or to
% \documentclass[reprint,aps,pra]{revtex4-2} for PRA.
% =============================================================================
\documentclass[11pt,letterpaper]{article}

\usepackage[margin=1in]{geometry}
\usepackage{graphicx, booktabs, amsmath, amssymb}
\usepackage{xr-hyper}
\usepackage[hidelinks]{hyperref}

\usepackage{natbib, xcolor, microtype, caption}
\captionsetup{font=small, labelfont=bf}

% \externaldocument must come AFTER natbib so that \citep/\citet/\citealp
% commands embedded in supplement.aux \@writefile entries (TOC/LOF/LOT) are
% defined when xr-hyper imports them.
\externaldocument[S]{supplement}

\graphicspath{
  {../step5_pretrain_origin_comparison/figures/}
  {../step6_unweighted_analysis/figures/}
  {../step6_integration_analysis/figures/}
  {../step6_pretrain_origin_comparison/figures/}
  {../step56_combined_synthesis/figures/}
  {../../notebooks/checkpoints_step5/unweighted/figures/}
  {../../notebooks/checkpoints_step5/integration/figures/}
  {../../notebooks/checkpoints_step6/unweighted/figures/}
  {../../notebooks/checkpoints_step6/integration/figures/}
}

\title{Pretrain-Loss-Weighting Choice for Neural-Network Density Functionals: \\
       A Cross-Step Comparison of Unweighted vs.\ Becke-Quadrature-Integrated Objectives}
\author{Alec P.\ Wills}
\date{\today}

\begin{document}
\maketitle

\begin{abstract}
We compare two pretrain-loss-weighting schemes for a neural-network exchange-correlation
functional based on the PBE generalized-gradient-approximation form
\citep{PerdewBurkeErnzerhof1996}: an unweighted per-grid-point mean-squared-error and an
integration-weighted MSE that multiplies each pointwise residual by the Becke quadrature
volume element \citep{Becke1988} and the LDA exchange-energy density. The comparison
spans two training-data regimes: ``step 5'' (8 architectures, 3 losses, 3 solvers,
training on H$_2$O alone, 72 specs per origin) and ``step 6'' (2 architectures, 5 losses,
3 solvers, 3 training-step budgets, training on H$_2$O+C$_2$H$_2$, 90 specs per origin).
Across all 10 architectures evaluated, integration-weighted pretraining produces a tighter
$F_x$ basin without exception (per-arch mean tightening 3.98$\times$ in step 5, ratio of
population means 4.25$\times$; per-arch range 3.0--4.7$\times$ in step 5 and
3.0$\times$/3.1$\times$ for the two step-6 architectures). The tightening compounds through
downstream training in step 5 (best AE-MAE improves 11.9$\times$, a 2.80$\times$
multiplier on the 4.25$\times$ pretrain ratio); in step 6 the AE improvement (2.86$\times$)
is comparable to but does not exceed the pretrain-$F_x$ ratio (3.06$\times$), indicating
that the broader step-6 training set saturates the pretrain-origin advantage at the
extreme tail. The integration weighting effect is asymmetric between channels: $F_x$ is
tightened uniformly across all 10 architectures (population mean 4.25$\times$ in step 5),
while $F_c$ is essentially tied at the population level (1.03$\times$) with 5 of 8
step-5 architectures \emph{worse} under integration. Step 6's $V_{xc}$-aware losses
(L3, L4) produce the only specs in either pretrain origin with sub-$2\!\times\!10^{-4}$
e/bohr$^3$ density-RMSE on H$_2$O --- a 3$\times$ density-floor improvement over step 5,
unattainable without explicit $V_{xc}$ supervision. No
pretrain-weighting variant closes the in-distribution / transfer generalization gap; the
bottleneck is training-set diversity, not pretrain-weighting choice. The Lieb--Oxford bound
\citep{LiebOxford1981} is enforced architecturally (PBE's stricter design choice
$F_x \le 1.804$) across all 180 step-6 specs (90 per origin); zero violations.
\end{abstract}

\section{Introduction}
\label{sec:intro}

The PBE exchange enhancement factor is defined through the integrand structure of the
exchange energy. Verbatim from \citet[eq.~10]{PerdewBurkeErnzerhof1996}:
\begin{equation}
\label{eq:Ex_integrand}
E_X = \int d^3 r\; n\, \varepsilon_X^{\mathrm{unif}}(n)\, F_X(s),
\qquad
\varepsilon_X^{\mathrm{unif}}(n) = -\tfrac{3 e^2 k_F}{4\pi},
\end{equation}
with $k_F = (3\pi^2 n)^{1/3}$ and reduced gradient $s = |\nabla n|/(2 k_F n)$.
PBE's eq.~10 retains the explicit $e^2$ from Gaussian units; the rest of the manuscript
takes $e^2 = 1$ in Hartree atomic units. The PBE analytic form for $F_x$
(their eq.~14) is
\begin{equation}
\label{eq:PBE_Fx}
F_X(s) = 1 + \kappa - \frac{\kappa}{1 + \mu s^2/\kappa},
\qquad \kappa = 0.804.
\end{equation}
PBE chooses $\kappa = 0.804$, giving $F_x(s\!\to\!\infty) = 1+\kappa = 1.804$, as a
\emph{stricter} sufficient condition for the Lieb--Oxford bound
\citep{LiebOxford1981}.\footnote{\citet{LiebOxford1981} prove
$E_{XC} \ge -C\,\int n^{4/3}\, d^3 r$ with $C = 1.68$
\citep[their Theorem~3.1, eq.~26]{LiebOxford1981}; \citet{PerdewBurkeErnzerhof1996}
\S III(g) eq.~13 quote a tighter $C = 1.679$, with the same Lieb--Oxford-1981
attribution. The 0.001 difference is a numerical refinement of the same proof,
not a different bound; both round to ``1.68''.}
The bare bound permits up to $F_x \le 2.273$; verbatim from
\citet{PerdewBurkeErnzerhof1996}, \S III(g):
\begin{quote}
``The Lieb-Oxford bound [14]
$E_X\{n_\uparrow,n_\downarrow\}\geq E_{XC}\{n_\uparrow,n_\downarrow\}\geq -1.679\, e^{2}\int d^3 r\, n^{4/3}$
will be satisfied if the spin-polarized enhancement factor
$F_X(\zeta=1, s) = 2^{1/3} F_X(s/2^{1/3})$ grows gradually with $s$ to a maximum value
less than or equal to 2.273, i.e., $F_X(s) \le 1.804$.''\footnote{The bracketed
``[14]'' in the PBE original cites \citet{LiebOxford1981}; the inequality is
labelled eq.~(13) in PBE 1996. The explicit $e^2$ factor is from the
Gaussian-units form of the Lieb--Oxford theorem; in the atomic-unit convention
used elsewhere in this manuscript $e^2 = 1$.}
\end{quote}
Our network (\texttt{xcquinox/alec/networks.py:\_AlecLOB}) enforces $F_x \le 1.804$
architecturally via a sigmoid clamp whose ceiling is an immutable static field; gradient
descent cannot relax it. The Lieb--Oxford theorem is therefore satisfied at every grid
point regardless of pretrain weighting, training loss, or architecture choice.

The integration-weighted pretrain loss is motivated by Eq.~\ref{eq:Ex_integrand} directly:
the integrand $\varepsilon_X^{\text{unif}}(n)\, F_x(s) \propto \rho^{4/3}\, F_x$ is biggest
at high-density bonding-region grid points, so pretrain-error reduction concentrated there
should produce a network whose $F_x$ is tighter exactly where it matters for downstream
energy predictions. We test this hypothesis empirically.

\section{Methods}
\label{sec:methods}

\subsection{Network architecture}
The $F_x$ network is an MLP gated by $\tanh^2(s)$, written as
$g(\boldsymbol{x},s) = \tanh^2(s)\,\mathrm{MLP}(\boldsymbol{x})$, then passed through a
sigmoid clamp
\begin{equation}
F_x = \mathrm{limit}\cdot \sigma\!\bigl(g - \log(\mathrm{limit}-1)\bigr),
\qquad \mathrm{limit} = 1.804,
\end{equation}
implemented in \texttt{xcquinox/alec/networks.py:42--44} as
\texttt{self.limit$\,*\,$jax.nn.sigmoid(x$\,-\,$jnp.log(self.limit$\,-\,$1))$\,-\,$1}
and wrapped at \texttt{networks.py:120--123} via \texttt{return 1$\,+\,$lobterm}.
The $-\log(\mathrm{limit}-1)$ bias makes the sigmoid centred so that $g=0$ yields
$F_x \approx 1$ (LDA limit); range $F_x \in [0, 1.804]$, with the PBE design ceiling
$1+\kappa = 1.804$ approached only as $g\to\infty$.
The attention variant adds a multi-head scaled-dot-product attention block
\citep{VaswaniAttention2017} (implemented as
\texttt{SelfAttentionBlock} at
\texttt{xcquinox/net.py:50--136}, with Q/K/V/O linear projections
\texttt{Linear(H,H)}, \texttt{num\_heads}-way reshape to head dim
\texttt{H/num\_heads}, scaled-dot-product softmax along the last
axis, output projection, and a residual+LayerNorm wrap), applied with
the Pre-LayerNorm convention
(\citet{XiongPreLN2020} provides the warmup-free theoretical analysis;
the Pre-LN architecture itself was proposed earlier in the literature)
after the first MLP layer. Verbatim from \citet{VaswaniAttention2017}
\S 3.2.1 (the attention equation that our \texttt{SelfAttentionBlock} implements):
\begin{equation*}
\mathrm{Attention}(Q, K, V) = \mathrm{softmax}\!\left(\frac{Q K^T}{\sqrt{d_k}}\right) V .
\end{equation*}
The architectural Lieb--Oxford clamp is shared across all 10 architectures; gradient
descent has no path to mutate the ceiling.

\subsection{Pretrain weighting}

The two pretrain runs minimize different objectives over the same dataset:
\begin{align}
\mathcal{L}_{\text{unweighted}} &= \tfrac{1}{N_g} \sum_{g=1}^{N_g} \big(F_x^{\text{NN}}(s_g) - F_x^{\text{PBE}}(s_g)\big)^2 \\
\mathcal{L}_{\text{integration}} &= \frac{\sum_{g=1}^{N_g} W_g\,
   \big(F_x^{\text{NN}}(s_g) - F_x^{\text{PBE}}(s_g)\big)^2}
   {\sum_{g=1}^{N_g} W_g + \epsilon},
\qquad W_g = w_g^{\text{grid}}\,
   \big|\rho_g\, \varepsilon_x^{\text{LDA}}(\rho_g)\big| .
\end{align}
Here $w_g^{\text{grid}}$ is the Becke partitioning weight at grid point $g$
\citep{Becke1988}, $|\rho \varepsilon_x^{\text{LDA}}|$ is the magnitude of the LDA
exchange-energy density, and $\epsilon = 10^{-12}$ is a numerical regulariser. The
loss is a \emph{weighted mean} (not a bare weighted sum), implemented at
\texttt{xcquinox/alec/pretrain.py:288--291} as
\texttt{jnp.sum(w$\,*\,$residual\_sq)$\,/\,$(jnp.sum(w)$\,+\,$1e-12)}; this keeps the
loss commensurate with the unweighted MSE in scale while concentrating gradient
pressure on the high-$|\rho \varepsilon_x|$ region where $E_X$'s integrand is largest.
The corresponding correlation-channel loss uses an analogous weight built from the
PW92 correlation energy density $\varepsilon_c^{\mathrm{PW92}}$ \citep{PerdewWang1992}
via \texttt{pw92c\_unpolarized\_scalar} (\texttt{xcquinox/alec/pretrain.py:69--77}),
so that the $F_c$ residual is weighted by its own integrand structure, not by LDA
exchange.

\subsection{Architectures, losses, and solvers}
\label{sec:methods_arch_loss}

Step 5 enumerates eight architectures, factoring three orthogonal
augmentations on a base \texttt{deep} MLP: optional density-matrix-overlap
descriptors (\texttt{\_dm}), optional cusp-correction features
(\texttt{\_cusp}), and optional combined density-matrix~$+$~cusp
descriptors (\texttt{\_combined}); each variant is then either with or
without a self-attention block (\texttt{\_attn}). The eight specs are
\texttt{deep}, \texttt{deep\_attn}, \texttt{deep\_dm},
\texttt{deep\_dm\_attn}, \texttt{deep\_cusp}, \texttt{deep\_cusp\_attn},
\texttt{deep\_combined}, \texttt{deep\_combined\_attn}. Step 6 retains
only the two descriptor-richest variants
(\texttt{deep\_combined} and \texttt{deep\_combined\_attn}).

Step 5's three task losses are \texttt{A\_atomization} (AE only),
\texttt{B\_atomization\_plus\_dm} (AE plus a density-matrix residual),
and \texttt{C\_atomization\_plus\_grid} (AE plus an integrated grid-density
residual). Step 6 replaces these with five $V_{xc}$-aware variants:
\texttt{L1\_B} (the step-5 \texttt{B} carried forward),
\texttt{L2\_C\_anchor} (\texttt{C} plus a closed-shell-asymptote anchor on
$F_x(s\!\to\!\infty)\!\to\!1.804$),
\texttt{L3\_balanced\_vxc} (AE plus an explicit $V_{xc}$ residual against
the analytic PBE potential),
\texttt{L4\_balanced\_vxc\_anchor} (\texttt{L3} plus the
$F_x$-asymptote anchor), and
\texttt{L5\_gradnorm\_vxc} (\texttt{L3} with GradNorm-balanced multitask
weights, \citealp{ChenGradNorm2018}, $\alpha=1.5$).

The three solvers are \texttt{oneshot} (single non-self-consistent pass
on the precomputed reference density), \texttt{fixed\_j\_3} (fixed-point
SCF with three Coulomb-integral re-evaluations), and \texttt{full\_3}
(three full self-consistent SCF cycles). Step~6 adds three training
schedules: \texttt{group1} (H$_2$O, short), \texttt{group2}
(H$_2$O+C$_2$H$_2$, short), \texttt{group3} (H$_2$O+C$_2$H$_2$, long).

\subsection{Training and evaluation}
\label{sec:methods_training_eval}

\paragraph{Network shape.} The MLP body has \texttt{depth}=4 hidden layers
of \texttt{nodes}=32 width with GELU activations
(\texttt{xcquinox/alec/config.py} entry \texttt{deep}; verified at
\texttt{xcquinox/alec/networks.py:84}); descriptor-augmented variants
(\texttt{\_dm}, \texttt{\_cusp}, \texttt{\_combined}) prepend the
corresponding extra features to the input vector. The attention variant
inserts the \texttt{SelfAttentionBlock} (\S\ref{sec:methods_arch_loss})
between the first and second hidden layers with \texttt{num\_heads}=4
(\texttt{config.py} entry \texttt{deep\_attn}).

\paragraph{Optimization.} Pretrain and task-loss phases share the same
optimizer chain (\texttt{xcquinox/alec/pretrain.py:\_build\_optimizer},
\texttt{xcquinox/alec/train.py:build\_optimizer}):
\texttt{optax.clip\_by\_global\_norm(1.0)} followed by
\texttt{optax.adam} with a linear-decay learning-rate schedule from
\texttt{lr\_start}=$10^{-2}$ to \texttt{lr\_end}=$10^{-5}$ over the run
length, with the decay starting at \texttt{lr\_decay\_start}=20\% of
total steps. No weight decay; no warmup.

\paragraph{Pretrain data and reference targets.} The pretrain set comprises
the four atomic species $\{H, He, O, N\}$ (singlet He plus open-shell
H/N/O reflecting first/second-row coverage relevant to step 5 and step 6
training molecules) plus the molecular pretrain probe (H$_2$O for both
steps; step 6 additionally probes C$_2$H$_2$ for analytic checks). The
pretrain target $F_x^{\mathrm{PBE}}(s)$ comes from the analytic PBE form
(eq.~\ref{eq:PBE_Fx}); $F_c^{\mathrm{PBE}}$ comes from the analytic
PBE correlation enhancement.

\paragraph{Training molecules and basis.} Step 5 trains on H$_2$O for 250
task-loss steps; step 6 trains on H$_2$O+C$_2$H$_2$ for 100 (groups~1,~2)
or 250 (group~3) task-loss steps. Geometries are taken from the W4-11
benchmark set (CCSD(T)-CBS-quality reference geometries; equilibrium
bond distances, no out-of-equilibrium sampling). The basis is
\texttt{def2-svp} throughout; the integration grid is PySCF
\texttt{grid\_level}=1, the Becke partitioned multicenter grid
\citep{Becke1988}.

\paragraph{Reference atomization energies.} AEs are evaluated against the
W4-11 reference values (140 first/second-row total atomization energies
derived from W4-quality \textit{ab initio} data, with $3\sigma$ uncertainty
under 1~kJ/mol relative to Active Thermochemical Tables;
\citealp{KartonDaonMartin2011}). Densities are integrated on the same
Becke--Lebedev grid that pretraining uses. Atomic correlation reference
totals for first/second-row atoms (C, N, O) follow
\citet{ChakravortyEtAl1993}, whose primary tabulation is 3-- through
18-electron species. The H atom is treated analytically
($E^{\mathrm{exact}}_{\mathrm{H}} = -\tfrac{1}{2}\,\mathrm{Ha}$, no
correlation by definition), and He uses an explicit non-Chakravorty
reference (Chakravorty's Table~VIII includes a re-analysis of 2-electron
data via Davidson, but is not its primary source).

\paragraph{Test and transfer sets.} ``In-distribution'' AE-MAE refers to
evaluation on the training molecule(s) themselves (H$_2$O for step 5;
H$_2$O+C$_2$H$_2$ for step 6, with results sliced per-group as group~1
through group~3). ``Transfer'' refers to molecules absent from the
training set: a primary transfer set $\{$CH$_4$, H$_2$, OH$\}$ shared by
both step subsets, and (step 6 only) a secondary transfer set
$\{$CO$_2$, HF, NH$_2$, NH$_3\}$ probing closed/open-shell light-hydride
generalisation. Cross-step transfer comparisons in this manuscript
restrict to the shared primary set.

\section{Results}
\label{sec:results}

\subsection{Combined headline}
\label{sec:results_headline}

Table~\ref{tab:headline_4way} summarizes the four-workflow comparison.
Within each step, the integration-pretrain origin produces a tighter best
AE-MAE on H$_2$O at the same training budget --- 11.9$\times$ tighter in
step 5 (0.001402 vs.\ 0.000118 kcal/mol) and 7.3$\times$ tighter in step 6
(0.003112 vs.\ 0.000426 kcal/mol on H$_2$O-restricted rows). Both step
subsets show the population-mean pretrain-$F_x$ ratio (4.25$\times$ in
step 5, 3.06$\times$ in step 6); the downstream best-AE
ratio is greater in step 5 (because narrow training amplifies the
pretrain head start) and saturates in step 6.
The integration-origin step-6 best $\rho$-RMSE
($1.23\!\times\!10^{-4}$ e/bohr$^3$) is the tightest density of any
spec in any of the four workflows.

\begin{table}[h]
\centering
\caption{Cross-workflow headline numbers. ``best AE on H$_2$O''
restricts step-6 evaluations to H$_2$O rows so the cross-step axis is
apples-to-apples.}
\label{tab:headline_4way}
\small
\begin{tabular}{lrrrr}
\toprule
workflow & $n$ & best AE H$_2$O & best $\rho$-RMSE H$_2$O & pretrain $F_x$ mean \\
& & (kcal/mol) & (e/bohr$^3$) & \\
\midrule
step5 / unweighted  & 72 & 0.001402 & $4.41\!\times\!10^{-4}$ & $1.50\!\times\!10^{-4}$ \\
step5 / integration & 72 & \textbf{0.000118} & $3.72\!\times\!10^{-4}$ & $3.52\!\times\!10^{-5}$ \\
step6 / unweighted  & 90 & 0.003112 & $1.32\!\times\!10^{-4}$ & $3.76\!\times\!10^{-5}$ \\
step6 / integration & 90 & 0.000426 & \textbf{$1.23\!\times\!10^{-4}$} & \textbf{$1.23\!\times\!10^{-5}$} \\
\bottomrule
\end{tabular}
\end{table}

\subsection{Per-architecture pretrain-$F_x$ comparison (Fig.~\ref{fig:comb2})}
\label{sec:results_perarch_fx}

The pretrain-$F_x$ MSE comparison (Fig.~\ref{fig:comb2}) is the most direct
test of the integrand-structure hypothesis: integration weighting concentrates
gradient pressure where $E_X$'s integrand is largest, so the pretrained $F_x$
should be tighter at exactly the points that matter most for the energy.
Across all 10 architectures, integration tightens the final pretrain $F_x$
without exception. In step 5 the per-arch ratios are
$\{4.74, 4.74, 4.06, 4.03, 3.82, 4.51, 2.95, 3.01\}$ for
\texttt{deep}, \texttt{deep\_attn}, \texttt{deep\_dm}, \texttt{deep\_dm\_attn},
\texttt{deep\_cusp}, \texttt{deep\_cusp\_attn}, \texttt{deep\_combined},
\texttt{deep\_combined\_attn} respectively
(per-arch mean $3.98\times$, range $2.95$--$4.74$;
verified against \texttt{notebooks/checkpoints\_step5/\{uw,int\}/pretrain/<arch>/pretrain\_metadata.json}).
The descriptor-rich \texttt{\_combined} variants show the smallest gain
($\sim 3\times$); the simpler \texttt{deep} and \texttt{deep\_attn} bases
show the largest ($\sim 4.7\times$), consistent with extra descriptor capacity
partly offsetting the unweighted-origin's wasted gradient on low-density
tail points. Step 6's two architectures give $3.00\times$ (\texttt{deep\_combined})
and $3.13\times$ (\texttt{deep\_combined\_attn}) --- consistent with their
step-5 counterparts ($2.95\times$ and $3.01\times$), so the per-arch ratio is
not strongly affected by the broader step-6 training set.

\begin{figure}[h]
\centering
\includegraphics[width=\textwidth]{figcomb_2_pretrain_fx_4way.png}
\caption{Pretrain final $F_x$ MSE per architecture (log scale).
(a)~Step 5's 8 architectures; (b)~step 6's 2 architectures, a subset of step 5's.
Integration weighting tightens $F_x$ on every architecture without exception;
the per-arch ratios range from $\sim$3.0 to $\sim$4.7. This is the cleanest
cross-step physical claim of the comparison.}
\label{fig:comb2}
\end{figure}

\subsection{Pretrain-$F_c$ comparison and the F$_x$/F$_c$ asymmetry}
\label{sec:results_fc}

The same integration weighting is applied independently to the correlation
channel via the PW92 correlation energy density
$|\rho\,\varepsilon_c^{\mathrm{PW92}}|$ \citep{PerdewWang1992}, but the
$F_c$ pretrain shows a much weaker integration advantage. The
population-mean ratio is $1.03\times$ (essentially tied:
$4.80\!\times\!10^{-5}$ unweighted vs.\ $4.66\!\times\!10^{-5}$ integration);
moreover, on a per-architecture basis 5 of 8 step-5 architectures are
\emph{worse} under integration on $F_c$ (\texttt{deep\_dm},
\texttt{deep\_dm\_attn}, \texttt{deep\_cusp\_attn}, \texttt{deep\_combined},
\texttt{deep\_combined\_attn}). The 3 that benefit
(\texttt{deep}, \texttt{deep\_attn}, \texttt{deep\_cusp}) include
\texttt{deep} and \texttt{deep\_attn}, the two architectures with the
largest absolute unweighted $F_c$ baselines ($1.15\!\times\!10^{-4}$ and
$9.99\!\times\!10^{-5}$, vs.\ $\sim 4 \!\times\!10^{-5}$ for the rest);
their per-arch ratios ($1.63\times$ and $1.30\times$ in favour of
integration) dominate the population mean despite the per-arch tally
running 5/8 against integration.

The asymmetry between $F_x$ (uniformly tightened) and $F_c$ (mixed)
under the same weighting strategy is consistent with two independent
factors. (a) PW92 $\varepsilon_c$ is roughly an order of magnitude smaller
in magnitude than $\varepsilon_x^{\mathrm{LDA}}$ at typical bonding densities,
so the $F_c$-weighted residual has a narrower dynamic range and a smaller
high-density-region gradient signal. (b) PBE's $F_c$ functional form is
already much closer to the LDA limit than $F_x$ is in the GGA regime
(modest gradient enhancement), leaving less ``room to move'' under any
re-weighting. We do not attempt to disentangle these contributions in this
manuscript; the F$_c$/F$_x$ asymmetry is documented as a robust empirical
observation across both step subsets.

\subsection{Pretrained-only NN baseline ladder}
\label{sec:results_pretrained_only}

Before downstream task training, we evaluate the network state immediately
after the 1000-step pretrain phase as a self-contained baseline. In step 5
the pretrained-only NN AE on H$_2$O is 27.2~kcal/mol (unweighted) vs.\
17.1~kcal/mol (integration) --- a 1.59$\times$ ratio, equivalent to a 37.1\%
AE-MAE reduction. In step 6 the pretrained-only AE-MAE on
H$_2$O+C$_2$H$_2$ is 121.4~kcal/mol (unweighted) vs.\ 119.0~kcal/mol
(integration) --- only 1.02$\times$, statistically tied. The shrinking
pretrain-origin advantage from step 5 to step 6 reflects step 6's broader
training set (H$_2$O+C$_2$H$_2$): with two molecules, the unweighted-origin
pretrain has access to enough density-region diversity at the start of
training that the integration head start is largely neutralised before
gradient descent even begins. The downstream task-loss optimization then
re-opens the gap only at the best-spec extreme tail
(Table~\ref{tab:headline_4way}); the population median is much closer
between origins.

\subsection{Density vs.\ energy plane (Fig.~\ref{fig:comb3})}
\label{sec:results_density_vs_energy}

The density-vs-energy plane (Fig.~\ref{fig:comb3}) places each spec at
$(\text{AE-MAE}, \rho\text{-RMSE})$ on H$_2$O and colour-codes by workflow.
Two qualitative observations dominate. First, step 6's V$_{xc}$-aware
losses (L3, L4) are the only specs in any of the four workflows that
reach sub-$2\!\times\!10^{-4}$ e/bohr$^3$ density-RMSE; this is a 3$\times$
improvement on step 5's best density floor ($3.72\!\times\!10^{-4}$),
unattainable without explicit $V_{xc}$ training. Second, within each step,
integration shifts the cluster left (lower AE) without raising the
density-RMSE ceiling --- a Pareto improvement on both axes. Across steps,
step 6 sits to the lower-right of step 5 (better density, worse AE on H$_2$O
because step 6 averages over a multi-molecule training set), confirming
that pretrain weighting and training-set diversity are independent axes.

\begin{figure}[h]
\centering
\includegraphics[width=\textwidth]{figcomb_3_density_vs_ae_unified.png}
\caption{Each spec is one point at $(\text{AE-MAE}, \rho\text{-RMSE})$ on H$_2$O.
Color encodes workflow. Step 6's V$_{xc}$-aware losses (L3, L4) are the only
specs in any workflow that reach sub-$2 \times 10^{-4}$ e/bohr$^3$ density-RMSE;
within each step, integration shifts the cluster left (lower AE) without
raising density-RMSE. The within-run Pareto observation is consistent with
Burke's earlier formulation of the energy-vs-density tradeoff (cf.~paraphrased
in \citet{Kepp2017}); the broader Medvedev across-decades thesis
\citep{Medvedev2017} is contested.}
\label{fig:comb3}
\end{figure}

\subsection{Best AE per workflow (Fig.~\ref{fig:comb1})}
\label{sec:results_bestae}

The best AE-MAE per workflow on H$_2$O (Fig.~\ref{fig:comb1}) is plotted
on a logarithmic axis. All four workflows clear chemical accuracy
\citep{Pople1999} on H$_2$O by orders of magnitude: 0.001402 kcal/mol
(step 5 unweighted, $5{,}014\times$ below PBE-on-H$_2$O), 0.000118 kcal/mol
(step 5 integration, $59{,}500\times$ below PBE), 0.003112 kcal/mol
(step 6 unweighted), and 0.000426 kcal/mol (step 6 integration). The
step-5 winner is \texttt{deep\_dm\_attn / C\_atomization\_plus\_grid /
fixed\_j\_3} on the integration origin; the step-6 winner is
\texttt{group3 / deep\_combined\_attn / L5\_gradnorm\_vxc / full\_3}
on the integration origin. In both step subsets the integration origin
takes the headline; the step-5 integration spec also achieves the tightest
in-distribution AE of any spec considered (a single-molecule, narrow-data
artefact, but a striking demonstration of how far gradient-descent can drive
the in-distribution residual when pretrain and task losses align).

\begin{figure}[h]
\centering
\includegraphics[width=\textwidth]{figcomb_1_best_ae_4way.png}
\caption{Best H$_2$O AE-MAE per workflow on a logarithmic axis,
with PBE-on-H$_2$O ($\sim 7.03$ kcal/mol; this run's specific def2-svp /
grid\_level=1 configuration) and chemical accuracy
(1 kcal/mol; \citet{Pople1999}) drawn as references.
The step-5 / integration winner achieves $1.18\times 10^{-4}$ kcal/mol,
$\sim$59,500$\times$ below PBE on the same H$_2$O configuration.}
\label{fig:comb1}
\end{figure}

\subsection{Distribution of H$_2$O AE-MAE across all specs (Fig.~\ref{fig:comb5})}
\label{sec:results_distribution}

The headline best-spec ratios in Table~\ref{tab:headline_4way} are extreme-tail
quantities; we therefore also examine the full population distribution
(Fig.~\ref{fig:comb5}, stacked histogram of $\log_{10}|AE\text{ error}|/$kcal/mol
per spec). Both step-5 distributions sit roughly two decades to the left of
both step-6 distributions, which is expected: step 5 trains on H$_2$O alone
and is evaluated on H$_2$O alone, so it is genuinely in-distribution at
training time. Within each step, the population-median shift between
unweighted and integration is small. Step 5's median moves $\sim$0.07~dex
leftward (integration tighter); step 6's median moves $\sim$0.4~dex
\emph{rightward} (integration slightly worse at the median). The tail
of the distribution behaves differently: integration's leftmost tail extends
roughly 1 decade further left than unweighted's in step 5, generating the
11.9$\times$ best-spec ratio in Table~\ref{tab:headline_4way}, but the
population median tells a much more modest story. The cross-step shift
($\sim$2 decades) dominates the cross-origin shift ($\lesssim$0.4 decades)
in both directions: training-set diversity matters more than pretrain
weighting at the population level.

\begin{figure}[h]
\centering
\includegraphics[width=\textwidth]{figcomb_5_normalized_landscape.png}
\caption{Stacked histogram of $\log_{10}(|AE \text{error}| / 1\,\text{kcal/mol})$
on H$_2$O across all specs in each workflow. Both step-5 distributions sit
roughly two decades left of both step-6 distributions. The within-step
unweighted-to-integration median shift is small in step 5
($\sim$0.07~dex leftward) but slightly \emph{rightward} in step 6
($\sim$0.4~dex worse on the H$_2$O population median, recomputed from
\texttt{run\_combined\_synthesis.py:plot\_normalized\_landscape}); the
across-step shift dominates, driven by step's narrower training
(H$_2$O alone in step 5) rather than by pretrain weighting alone. The
best-spec extreme tail shifts visible in Table~\ref{tab:headline_4way}
are not reflected at the population median.}
\label{fig:comb5}
\end{figure}

\subsection{Transfer-set comparison (Fig.~\ref{fig:comb4})}
\label{sec:results_transfer}

The transfer-set comparison (Fig.~\ref{fig:comb4}) restricts to the three
molecules common to step 5's and step 6's transfer sets ($\{$CH$_4$, H$_2$,
OH$\}$) and to the cross-step-analogous loss families: $\{B,\,C\}$ in
step 5 vs.\ $\{$L1\_B, L2\_C\_anchor$\}$ in step 6. Step 6's exclusively
$V_{xc}$-aware losses (L3, L4, L5) have no step-5 analog and are excluded
to avoid an apples-to-oranges median.

Three patterns emerge. CH$_4$ favours step 5 (median 5--7~kcal/mol vs.\
$\sim$10--15 for step 6), reflecting that CH$_4$'s closed-shell, sp$^3$,
single-bonded character is closer to H$_2$O (which trains step 5
exclusively) than to C$_2$H$_2$ (a triple-bonded, $\pi$-rich molecule
that broadens step 6's training set away from the CH$_4$ region).
H$_2$ marginally favours step 6, where the broader training set provides
some bond-stretching diversity absent from H$_2$O alone. OH is uniformly
above PBE in all four workflows (5--21~kcal/mol vs.\ PBE's 1.84~kcal/mol)
because OH is a $^2\Pi$ open-shell radical and the molecular pretrain
set $\{$H$_2$O$\}$ (step 5) or $\{$H$_2$O, C$_2$H$_2\}$ (step 6) has no
open-shell analog; the atomic pretrain set $\{H, He, O, N\}$ has open-shell
atoms but cannot supply the molecular $\pi$-orbital spin-density profile
that determines OH's exchange-correlation behaviour. Pretrain-weighting
moves the median by $<\!2\times$ on every transfer molecule --- much
smaller than the across-workflow factor-of-3-or-more shifts, confirming
that the bottleneck for transfer is training-set composition, not the
choice of weighting scheme.

\begin{figure}[h]
\centering
\includegraphics[width=\textwidth]{figcomb_4_transfer_overlap.png}
\caption{Restricted-loss-set transfer comparison. Step 5's losses
$\{B,C\}$ are compared against step 6's $\{$L1\_B, L2\_C\_anchor$\}$ (the
direct cross-step analogs); step 6's V$_{xc}$-aware losses (L3, L4, L5)
have no step-5 analog and are excluded to avoid inflating step 6's
median. Per-mol PBE references: CH$_4$ 0.91, H$_2$ 7.31, OH 1.84
kcal/mol. CH$_4$ favors step 5's narrower training; H$_2$ slightly
favors step 6's broader set; OH is uniformly above PBE across all
four workflows due to its open-shell ($^2\Pi$) character not being
covered by the closed-shell molecular pretrain set.}
\label{fig:comb4}
\end{figure}

\subsection{Lieb--Oxford / PBE-design ceiling enforcement audit}
\label{sec:results_lob}

The architectural \texttt{\_AlecLOB} clamp guarantees
$F_x \in [0, 1.804]$ at every grid point regardless of training history.
We verify this empirically by post-processing every step-6 checkpoint:
the audit script
(\texttt{notebooks/analysis/audit\_lob\_enforcement.py}) loads each spec
on a CH$_4$ molecular grid, evaluates $F_x(s_g)$ at every grid point,
and counts the global maximum.

Across the 90 step-6 specs in the unweighted origin and the 90 in the
integration origin (180 total checkpoints), the global maximum
$F_x$ value observed is exactly 1.804000 to six decimals --- the design
ceiling, never exceeded. Zero specs in either origin produce any
$F_x > 1.804$ at any grid point. The minimum-$F_x$ value observed
(0.0335 in integration, 0.000 in unweighted) hits the lower sigmoid floor
in low-density tail regions, which is permitted by the clamp form
($F_x \to 0$ as $g \to -\infty$). The \texttt{\_AlecLOB.limit} field is
declared \texttt{eqx.field(static=True)} at
\texttt{xcquinox/alec/networks.py:21}, so gradient descent has no path
to mutate the ceiling --- the bound is enforced at the level of the
computation graph, not by post-hoc projection.

Recall (\S\ref{sec:intro}) that 1.804 is PBE's stricter design choice;
the bare Lieb--Oxford bound permits $F_x \le 2.273$
\citep{LiebOxford1981, PerdewBurkeErnzerhof1996}. A network with a
relaxed clamp at 2.273 would have additional fitting headroom but would
also be a less constrained approximation of the exact functional than
PBE itself. We adopt 1.804 throughout for direct comparability with PBE
references.

\section{Compelling conclusions}
\label{sec:conclusions}

\paragraph{(1) Integration weighting tightens $F_x$ pretrain quality on every architecture.}
Across all 10 architectures (8 in step 5, 2 in step 6), integration-weighted pretraining
produces a tighter final $F_x$ MSE in every single case. Per-arch ratio range
3.0--4.7$\times$ in step 5 (mean 3.98$\times$ per-arch, 4.25$\times$ for the population
means); the two step-6 architectures give 3.00$\times$ and 3.13$\times$ (population-mean
ratio 3.06$\times$). This is the cleanest cross-workflow finding of the synthesis
(Fig.~\ref{fig:comb2}). The integrand-structure justification follows directly from
\citet{PerdewBurkeErnzerhof1996} eq.~10 (Eq.~\ref{eq:Ex_integrand} in this manuscript):
pretrain pressure concentrated where $E_X$'s integrand is biggest produces a network whose
$F_x$ is tighter exactly there.

\paragraph{(2) The pretrain head start compounds through step-5 training but saturates in step 6.}
In step 5 the downstream best-AE improvement (11.9$\times$) is 2.80$\times$ the pretrain
$F_x$ improvement (4.25$\times$); the head start \emph{compounds} through 250 task-loss
steps. In step 6 the downstream best-AE improvement (2.86$\times$) is comparable to but
does not exceed the pretrain $F_x$ ratio (3.06$\times$); the broader training set
(H$_2$O+C$_2$H$_2$) and richer V$_{xc}$-aware losses already provide enough gradient signal
that the unweighted-origin run catches up at the best-spec extreme. There is no universal
``compounding factor'' across the two step subsets; step 5's compounding is a feature of
the H$_2$O-only narrow-data regime.

\paragraph{(3) V$_{xc}$-aware losses are necessary for sub-$2\times 10^{-4}$ density quality.}
Step 5's loss set lacks V$_{xc}$ fitting; step 5's best density-RMSE is $3.72\times 10^{-4}$
e/bohr$^3$. Step 6 added V$_{xc}$-LossNorm losses (L3, L4); their best density-RMSE is
$1.23\times 10^{-4}$ — a 3$\times$ improvement only achievable with explicit V$_{xc}$
training data. The within-run Pareto pattern (V$_{xc}$-aware specs lose AE accuracy in
exchange for density quality, Fig.~\ref{fig:comb3}) is consistent with Burke's earlier
formulation of the energy-vs-density tradeoff (paraphrased in \citet{Kepp2017}); the
broader Medvedev across-decades thesis \citep{Medvedev2017} is contested by
\citet{Kepp2017}, who shows the historical trend is dominated by one specific group's
recent functionals rather than a general DFT-development bias.

\paragraph{(4) Transfer is dominated by molecule-specific physics, not pretrain weighting.}
Across the three molecules common to step 5's and step 6's transfer sets — CH$_4$, H$_2$,
OH — pretrain weighting moves the median MAE by less than $2\times$. CH$_4$ favors step
5 (narrower training; bonding pattern closer to H$_2$O than to C$_2$H$_2$); H$_2$ slightly
favors step 6 (broader training); OH is uniformly above PBE in all four workflows because
it is open-shell ($^2\Pi$, $S=1/2$) while the molecular pretrain set is closed-shell.
This is the in-distribution-vs-holdout caution of \citet{BehlerParrinello2007} applied to
NN XC functionals (the original paper is for NN potential-energy surfaces; the
generalization principle applies but the citation is by analogy).

\paragraph{(5) F$_c$ pretrain is mostly insensitive to the integration weighting.}
The implementation already builds the F$_c$ weight from $|\rho\,\varepsilon_c^{\text{PW92}}|$
(\texttt{xcquinox/alec/pretrain.py:69--77}, using the LSDA correlation parameterization of
\citealp{PerdewWang1992}) — not from $|\rho\,\varepsilon_x^{\text{LDA}}|$. The F$_c$
stagnation under integration weighting therefore does not have a ``wrong prefactor'' cause;
likely candidates are (a) PW92 $\varepsilon_c$ being $\sim 10\times$ smaller in magnitude
than $\varepsilon_x^{\text{LDA}}$ at typical bonding densities (so the F$_c$-weighted
residual has a narrower dynamic range), or (b) F$_c$'s gradient enhancement intrinsically
having less room to move than F$_x$'s.

\paragraph{(6) Lieb--Oxford / PBE-design ceiling holds architecturally across all four
workflows.} The \texttt{\_AlecLOB} clamp is pretrain-origin-independent;
audit (\texttt{notebooks/analysis/audit\_lob\_enforcement.py}) on the 90 step-6
checkpoints in each origin shows $F_x \le 1.804$ in 90/90 specs. The ceiling 1.804 is
PBE's stricter design (\S\ref{sec:intro}); the bare Lieb--Oxford bound permits
$F_x \le 2.273$ \citep{LiebOxford1981, PerdewBurkeErnzerhof1996}.

\section{Limitations and honesty disclosures}
\label{sec:limits}

This section enumerates the places where conclusions rest on empirical observation
rather than first-principles derivation, or where citations are by analogy rather than
direct.

\begin{enumerate}
\item Cross-step $F_x$ mean ratios are not on equal footing — step 5 averages over 8
architectures (light to heavy), step 6 over 2 (both heavy). The fair cross-step
comparison is the per-arch overlay (Fig.~\ref{fig:comb2}), not the population means.

\item PBE-on-H$_2$O $\approx 7.03$ kcal/mol is run-specific to our def2-svp /
grid\_level=1 / W4-11-geometry configuration; it is \emph{not} a literature value.
W4-11 itself \citep{KartonDaonMartin2011} tabulates CCSD(T)-quality reference
atomization energies, not PBE errors.

\item Best-spec-per-workflow numbers are minimum AE-MAE values; a single tight spec
dominates the headline. The distribution histogram (Fig.~\ref{fig:comb5}) confirms
the cross-step difference is a population shift.

\item The compounding factor $\sim 2.8\times$ in conclusion (2) is empirical.

\item The transfer-set comparison combines step 5's transfer-mol output with step 6's
\texttt{transfer\_primary\_df} — these are not bit-identical pipelines (step 6 uses
different precompute paths and a CCSD-target overlay). Cross-step transfer numbers are
qualitatively but not quantitatively identical-protocol.

\item OH transfer: the pretrain atom set $\{H, He, O, N\}$ contains open-shell atoms
(only He has $S=0$), so the OH mismatch is on \emph{molecular} spin-density profile, not
on atomic open-shell character. An earlier draft of this report mis-stated this.

\item Step 6 evaluates two transfer sets: a primary set
$\{$CH$_4$, H$_2$, OH$\}$ (the same 3 mols as step 5) and a secondary set
$\{$CO$_2$, HF, NH$_2$, NH$_3\}$ (4 closed/open-shell light hydrides
absent from step 5). Cross-step transfer comparisons in this manuscript
restrict to the shared primary set; the secondary set is shown in
the step-6 supplements as additional generalisation diagnostics. NH$_3$
is in the step-6 secondary set, not the primary.

\item The Behler--Parrinello transferability framework \citep{BehlerParrinello2007}
is for NN potential-energy surfaces (their test system is bulk silicon
across the solid-semiconducting / liquid-metallic / high-pressure phases),
not NN XC functionals; the ``in-distribution MSE alone does not certify
out-of-distribution accuracy'' principle generalises to NN XC, but the
specific paper does not address NN XC.

\item The ``functionals which yield highly accurate energies often produce
potentials which differ markedly from the exact ones'' phrasing is paraphrased in
\citet{Kepp2017} and traces to \citet{Burke1998} (Chem.\ Phys.\ Lett.\ 265, 115;
the bib-key suffix is approximate, the published year is 1997); the underlying
Burke--Perdew--Ernzerhof work of the late 1990s
(adiabatic-connection method, hybrid functionals, etc.) is the canonical primary source
for the energy-vs-density tradeoff principle. The Medvedev 2017 across-decades thesis
\citep{Medvedev2017} is contested by \citet{Kepp2017}; our within-run finding does not
depend on Medvedev's stronger generalization.

\item The GradNorm asymmetry parameter $\alpha = 1.5$
\citep{ChenGradNorm2018} used by step 6's L5 spec is the optimum found
in the NYUv2+kpts (depth + surface normals + room-keypoints) experiment
of \citet[Fig.~7 and \S 7.2]{ChenGradNorm2018}; the same paper also
reports $\alpha = 0.12$ as optimal for a synthetic 2-task regression
toy problem (\S 5.4), so the value is task-specific and not asserted
as a universal default. For our DFT loss family (AE +
$\rho$-RMSE + V$_{xc}$ residual + anchor) the optimal $\alpha$ is
empirically reasonable but un-tuned; a controlled sweep is warranted.
\end{enumerate}

\bibliographystyle{plainnat}
\bibliography{references}

\end{document}
